\documentclass[12pt]{article}
\usepackage[margin=0.75in]{geometry}

\newcommand{\duedate}{03/25/2026}
\newcommand{\assignment}{Hopfield Networks is All You Need}

\newcommand{\name}{Aksel Kretsinger-Walters, adk2164}
\newcommand{\email}{adk2164@columbia.edu}

\makeatletter
\def\input@path{{../}{../../}{../../../}}
\makeatother
\input{pset_template_CLMM.tex} %% DO NOT CHANGE THIS LINE

\begin{document}
\psetheader %% DO NOT CHANGE THIS LINE
\section{Hopfield Networks is All You Need (Ramsauer et al., 2020)}

\subsection{Motivation}

Classical Hopfield networks are associative memory models that store patterns and retrieve them via energy
minimization. However, traditional formulations suffer from:
\begin{itemize}
		\item Limited storage capacity
		\item Shallow energy landscapes
		\item Difficulty scaling to modern machine learning settings
\end{itemize}

This paper revisits Hopfield networks and shows that a modernized formulation can scale to high-dimensional
settings and is closely related to attention mechanisms.

\subsection{Key Idea}

The paper introduces \textbf{modern Hopfield networks} with continuous states and an updated energy function.
The central insight is:

\begin{quote}
		Modern Hopfield updates are mathematically equivalent to attention mechanisms.
\end{quote}

This connects associative memory with transformer architectures.

\subsection{Model Formulation}

Given stored patterns $\{x_i\}_{i=1}^N$, the network defines an energy function:

\[
		E(\xi) = -\frac{1}{\beta} \log \left( \sum_{i=1}^N \exp(\beta x_i^T \xi) \right)
\]

The update rule becomes:

\[
		\xi \leftarrow \sum_{i=1}^N \text{softmax}(\beta x_i^T \xi) \; x_i
\]

This is equivalent to a weighted sum of stored patterns using a softmax over similarities.

\subsection{Connection to Attention}

This update rule is identical to attention:

\[
		\text{Attention}(q, K, V) = \text{softmax}(q K^T) V
\]

Mapping:
\begin{itemize}
		\item Query $q \leftrightarrow \xi$
		\item Keys $K \leftrightarrow x_i$
		\item Values $V \leftrightarrow x_i$
\end{itemize}

Thus, attention can be interpreted as a retrieval step in a Hopfield network.

\subsection{Memory Properties}

Modern Hopfield networks have dramatically improved capacity:

\begin{itemize}
		\item Exponential storage capacity in dimension
		\item Ability to store continuous patterns
		\item Fast convergence (often one update step)
\end{itemize}

They can retrieve:
\begin{itemize}
		\item Exact stored patterns
		\item Averaged patterns
		\item Global prototypes
\end{itemize}

\subsection{Energy Landscape}

The energy function induces different types of fixed points:

\begin{itemize}
		\item Individual memories (sharp minima)
		\item Metastable states (clusters)
		\item Global averages (broad minima)
\end{itemize}

This explains why attention can interpolate between patterns.

\subsection{Applications}

The paper demonstrates that modern Hopfield layers can be used in:

\begin{itemize}
		\item Multiple instance learning
		\item Set-based learning tasks
		\item Deep neural networks as a drop-in module
\end{itemize}

They provide a principled interpretation of attention as memory retrieval.

\subsection{Key Contributions}

\begin{itemize}
		\item Reformulation of Hopfield networks with continuous states
		\item Proof of equivalence between Hopfield updates and attention
		\item Exponential increase in storage capacity
		\item Integration of associative memory into deep learning
\end{itemize}

\subsection{Takeaway}

The paper shows that:

\[
		\text{Attention} \approx \text{Modern Hopfield Retrieval}
\]

This unifies classical associative memory with modern transformer architectures and provides a theoretical
foundation for attention mechanisms.
\end{document}
